\section*{ID7598: Why Gravity Is Weak (Scaling Argument): A(r)}
\paragraph{Metadata.}
PiID: 6291642193994 | RTEID: RTE:P29504.7213:T4.6324 | UUID: 93d64a3e-fe70-5589-a330-80fd4ee48adb

\paragraph{Canonical Statement.}
Quantitatively, if single-packet acceleration contribution at radius \textbackslash{}(r\textbackslash{}) is \textbackslash{}(a\_0(r)\textbackslash{}), then total acceleration is \textbackslash{}(a(r)=N a\_0(r)\textbackslash{}). If \textbackslash{}(a\_0\textbackslash{}sim 10\textasciicircum{}\{-50\}\textbackslash{},\textbackslash{}text\{m/s\}\textasciicircum{}2\textbackslash{}) and \textbackslash{}(N\textbackslash{}sim 10\textasciicircum{}\{40\}\textbackslash{}), then you get macroscopic \textbackslash{}(a\textbackslash{}sim10\textasciicircum{}\{-10\}\textbackslash{},\textbackslash{}text\{m/s\}\textasciicircum{}2\textbackslash{}). This is why gravity is many orders smaller than EM: each EM carrier (photon, charge) interacts strongly per quantum; these vacuum packets are assumed to have tiny per-packet effect.

\paragraph{Canonical Equation.}
\[
a(r)=N a_0(r)
\]

\paragraph{Dictionary.}
Quantitatively, if single-packet acceleration contribution at radius \textbackslash{}(r\textbackslash{}) is \textbackslash{}(a\_0(r)\textbackslash{}), then total acceleration is \textbackslash{}(a(r)=N a\_0(r)\textbackslash{}) [ID7598]

\paragraph{Encyclopedia.}
Why Gravity Is Weak (Scaling Argument): A(r) is an algebraic definition, written a(r)=N a\_0(r). It is an algebraic definition expressing the quantity directly in terms of the variables on its right-hand side. It is built from N, a\_0, r, defining a(r). Within the Trawin Zero Point registry it serves as one element of the framework's mathematical narrative.
